% Authoritative LaTeX; edit this file, not the historical Markdown.
\chapter{The Smallest Possible Language Model}\label{chapter-3-the-smallest-possible-language-model}\label{app:smallest-model}
\section{A token ID is still not an
answer}\label{a-token-id-is-still-not-an-answer}

\Appref{tokens} ended at a precise boundary. Human text became UTF-8 bytes. A
model-specific tokenizer turned those bytes into token IDs. The request
owned the resulting sequence, and generated IDs could travel back
through token bytes and a strict UTF-8 framer to become streamed text.

\Chref{request-lifecycle} treated the numerical function as opaque, but recorded its output
for the prompt \texttt{I like}. We will now derive that output from the
actual 28-parameter fixture. Nothing in the model consults an answer table:
the four scores follow from a row lookup, a matrix-vector product and a bias.

What can a model actually do with a token ID?

Suppose the tokenizer emits \texttt{TokenId(2)} for \texttt{like}. The
integer \texttt{2} does not mean ``moderately positive.'' It is not
twice as semantic as token 1. It does not contain a direction toward
\texttt{Rust}, a probability, or a grammatical role. It is an index in
one model vocabulary. The next boundary must turn that discrete identity
into numerical data on which the model can compute.

This chapter builds that boundary with the smallest useful numerical
language model we can inspect completely:

\begin{quote}
\begin{enumerate}
\item text  \(\longrightarrow\)  token IDs  \(\longrightarrow\)  embedding lookup  \(\longrightarrow\)  hidden vector
\item \(\longrightarrow\)  output projection  \(\longrightarrow\)  vocabulary logits
\end{enumerate}
\end{quote}

There is no neural-network framework, general tensor library, BLAS call,
Transformer, attention mechanism, training loop, or model file. The
parameters are 28 visible \texttt{f32} values. The kernel is two nested
Rust loops. We will calculate every multiplication independently and
compare the whole result.

The model will produce \textbf{logits}: one unnormalized score for every
vocabulary token. It will not produce probabilities. A temporary
deterministic argmax selector keeps the existing end-to-end lifecycle
executable, but selection is outside the model. \Appref{sampling} will ask how
logits become a distribution and how an engine selects and feeds back
the next token. We stop at correct logits.

\begin{quote}
\textbf{FIRST PRINCIPLE} A token ID is a categorical identity, not a
meaningful scalar quantity. Numerical model semantics begin when that
identity selects learned parameters.
\end{quote}

\section{What makes this a language
model?}\label{what-makes-this-a-language-model}

A language model assigns scores or probabilities to linguistic
sequences. For next-token inference, we ask for a distribution over the
next vocabulary token given an observed history. For a history
\((x_0,x_1,\ldots,x_t)\), the target is

\begin{equation}
P\!\left(x_{t+1}\mid x_0,x_1,\ldots,x_t\right).
\end{equation}

{\def\LTcaptype{none} % do not increment counter
\begin{longtable}[]{@{}
  >{\raggedright\arraybackslash}p{(\linewidth - 2\tabcolsep) * \real{0.5000}}
  >{\raggedright\arraybackslash}p{(\linewidth - 2\tabcolsep) * \real{0.5000}}@{}}
\toprule\noalign{}
\begin{minipage}[b]{\linewidth}\raggedright
Symbol
\end{minipage} & \begin{minipage}[b]{\linewidth}\raggedright
Meaning
\end{minipage} \\
\midrule\noalign{}
\endhead
\bottomrule\noalign{}
\endlastfoot
\(V_{\mathrm{vocab}}\) & number of token identities in the bound
vocabulary \\
\(D\) & hidden-vector dimension \\
\(x_t\) & token identity at sequence position \(t\) \\
\(\mathbf{E}\) & embedding matrix owned by the model \\
\(\mathbf{W},\mathbf{b}\) & output projection owned by the model \\
\(\mathbf{h}\) & one forward call's hidden activation \\
\(\mathbf{z}\) & one forward call's logit vector \\
\end{longtable}
}

The output space matters: the possible outcomes are all
\(V_{\mathrm{vocab}}\) entries in the model vocabulary. The model must
therefore produce \(V_{\mathrm{vocab}}\) scores before a policy can
choose one token.

Our first numerical model makes a severe approximation:

\begin{equation}
P\!\left(x_{t+1}\mid x_0,\ldots,x_t\right)
\approx P\!\left(x_{t+1}\mid x_t\right)
\quad\text{(ENGINE-1 approximation).}
\end{equation}

It accepts a full sequence but uses only \texttt{x\_t}, the last token.
This is a first-order, neural bigram-like language model. It is still a
language model: its input represents token history, its task is
next-token prediction, and its output has one position per vocabulary
token. It is not a Transformer and cannot represent long-range context.

That simplicity is an advantage for the present question. If we began
with attention, normalization, residual connections, positional
encoding, dozens of layers, and billions of parameters, the first
numerical boundary would disappear inside machinery we could not yet
audit. Here, every output score will have a short arithmetic trail.

Neural language models predate Transformers. Bengio, Ducharme, Vincent,
and Jauvin described a neural probabilistic language model that learned
distributed representations of words along with a probability function
over sequences in 2003. Their model uses richer context and training
machinery than ENGINE-1. The relevant historical point is narrow: a
language model need not be a Transformer, and discrete vocabulary
identities can enter computation through learned vector representations.

\section{From a discrete ID to a
vector}\label{from-a-discrete-id-to-a-vector}

Let the vocabulary contain \(V_{\mathrm{vocab}}\) token identities and
let each identity have \(D\) components. Store those vectors as rows of
an \textbf{embedding matrix}:

\begin{equation}
\mathbf{E}\in\mathbb{R}^{V_{\mathrm{vocab}}\times D}.
\end{equation}

\(V_{\mathrm{vocab}}\) is the number of rows. \(D\), the \textbf{hidden
dimension}, is the number of values in each row. ENGINE-1 stores every
value as \texttt{f32} in contiguous row-major order.

If the input token is \(x\), the embedding operation is the row gather

\begin{equation}
\mathbf{h}=\mathbf{E}_{x,:},
\qquad x\in\{0,\ldots,V_{\mathrm{vocab}}-1\},
\quad \mathbf{h}\in\mathbb{R}^{D}.
\end{equation}

\(x\) is a scalar \texttt{TokenId}; \(\mathbf{h}\) is a vector of shape
\texttt{{[}D{]}}. The
\href{https://github.com/hermonai/inside-the-llm-engine/blob/astra-visual-rewrite/diagrams/model/embedding-row-lookup.txt}{embedding
lookup diagram} makes the access concrete:

\EngineFigure{embedding-gather}{A categorical ID selects one row and copies its values into a new activation. The row's position is not its meaning.}{fig:embedding-gather}

For \(x=2\), the result is

\begin{equation}
\mathbf{h}=[1.0,-0.5,2.0].
\end{equation}

This is a row selection, sometimes called a gather. ENGINE-1 computes
the row start at the element offset

\begin{equation}
o_{\mathrm{row}}=xD\quad[\mathrm{elements}].
\end{equation}

and copies the next \texttt{D} values. The scalar value 2 is used only
in address calculation. We do not multiply the vector by 2.

\subsection{Why not feed the integer directly into
arithmetic?}\label{why-not-feed-the-integer-directly-into-arithmetic}

Imagine that vocabulary revision A assigns:

\begin{quote}
\begin{enumerate}
\item 2  \(\longrightarrow\)  like
\item 3  \(\longrightarrow\)  Rust
\end{enumerate}
\end{quote}

and revision B assigns:

\begin{quote}
\begin{enumerate}
\item 2  \(\longrightarrow\)  Rust
\item 3  \(\longrightarrow\)  like
\end{enumerate}
\end{quote}

The language meanings swapped while the integers did not. Any rule that
treats 3 as inherently larger or closer to something than 2 would change
meaning when the vocabulary was renumbered. Embedding rows attach
numerical parameters to identities explicitly.

An embedding does not guarantee that each dimension has a clean human
label. We must not say that dimension 0 is ``programming,'' dimension 1
is ``affection,'' or dimension 2 is ``grammar'' merely because a tiny
hand fixture makes a story tempting. In a trained model, useful behavior
can be distributed across many dimensions. For ENGINE-1, \texttt{h} is
exactly this: the numerical representation of the current input token
consumed by the next computation.

\subsection{Is lookup secretly matrix
multiplication?}\label{is-lookup-secretly-matrix-multiplication}

We could represent token 2 as a one-hot vector:

\begin{verbatim}
[0, 0, 1, 0]
\end{verbatim}

and multiply it by \texttt{E}. The result would select row 2. That
equivalence can be useful mathematically, but it is a poor description
of ENGINE-1's execution. Materializing \texttt{V} values, almost all
zero, and multiplying through every row would hide the actual access
pattern. The code reads one contiguous row.

\begin{quote}
\textbf{FIRST PRINCIPLE} An embedding lookup reads the row named by a
token ID. It does not assign meaning to the ID's integer magnitude, and
it need not execute as dense matrix multiplication.
\end{quote}

\section{One score for every possible next
token}\label{one-score-for-every-possible-next-token}

The hidden vector represents the input. The model now needs one score
for each candidate next token. The result has \(V_{\mathrm{vocab}}\)
components.

Introduce an \textbf{output projection} matrix \(\mathbf{W}\) and bias
vector \(\mathbf{b}\):

\begin{equation}
\mathbf{W}\in\mathbb{R}^{V_{\mathrm{vocab}}\times D},
\qquad \mathbf{b}\in\mathbb{R}^{V_{\mathrm{vocab}}}.
\end{equation}

Each row \(\mathbf{W}_{i,:}\) belongs to candidate token \(i\). Compute

\begin{equation}
\mathbf{z}=\mathbf{W}\mathbf{h}+\mathbf{b},
\qquad \mathbf{z}\in\mathbb{R}^{V_{\mathrm{vocab}}}.
\end{equation}

For one candidate \(i\):

\begin{equation}
z_i=b_i+\sum_{j=0}^{D-1}W_{ij}h_j,
\qquad 0\le i<V_{\mathrm{vocab}}.
\end{equation}

This operation takes the \textbf{dot product} of row \texttt{i} and
\texttt{h}, then adds one bias. The
\href{https://github.com/hermonai/inside-the-llm-engine/blob/astra-visual-rewrite/diagrams/model/one-logit-dot-product.txt}{one-logit
diagram} connects vector notation to scalar arithmetic:

\EngineFigure{logit-reduction}{Each weight and activation enter multiplication independently. The three products and one bias determine the logit for ID 3.}{fig:logit-reduction}

Repeat that operation for all \(V_{\mathrm{vocab}}\) rows. Matrix
notation compresses the idea. The following reference pseudocode exposes
the scalar work; the current engine uses the checked GEMV layer introduced
in \Appref{matmul}, with shape checks and a separately applied bias:

\begin{verbatim}
for output in 0..vocab_size {
    let mut accumulator = output_bias[output];
    let row_start = output * hidden_dim;
    for dimension in 0..hidden_dim {
        accumulator +=
            output_weight[row_start + dimension] * hidden[dimension];
    }
    logits[output] = accumulator;
}
\end{verbatim}

The inner loop calculates one candidate score. The outer loop repeats it
for the vocabulary. There is no answer table in this code: changing an
embedding, projection weight, or bias changes the arithmetic that
produces the vector. The canonical
\href{https://github.com/hermonai/inside-the-llm-engine/blob/astra-visual-rewrite/diagrams/model/engine-1-overview.txt}{ENGINE-1
overview} places those immutable parameters beside request-owned history
and sampler state.

\subsection{``Linear'' layers usually include
bias}\label{linear-layers-usually-include-bias}

Strictly, \texttt{W\ h} is a linear map and \texttt{W\ h\ +\ b} is
affine when \texttt{b} is nonzero. Machine-learning libraries commonly
call the combined operation a linear layer. We will use \textbf{output
projection} for the role and write the bias explicitly so neither
terminology nor code can hide it.

\section{What a logit is---and is
not}\label{what-a-logit-isand-is-not}
\index{logits}

Each \texttt{z\_i} is a \textbf{logit}: an unnormalized numerical score
for candidate token \texttt{i}. ENGINE-1's complete result for
\texttt{like} will be:

\begin{verbatim}
token       logit
<eos>       -0.7
I            0.1
like         0.4
Rust         2.2
\end{verbatim}

The values do not sum to 1. They can be negative. A logit of 2.2 is not
a probability of 2.2, 220 percent, or ``2.2 times likely.'' A difference
of 1.8 is not directly a probability difference. We have not yet defined
the transformation needed for those interpretations.

For a deterministic argmax policy, the larger finite score wins, so
\texttt{Rust} would be selected here. That statement is about ordering
only. It does not make the largest score probability 1, and it does not
make argmax part of the model.

The boundary remains:

\begin{verbatim}
Model::forward(history) -> ForwardPass { logits, ... }
SamplerState::sample(logits) -> SamplingStep { token_id, ... }
\end{verbatim}

The historical ENGINE-0 prototype blurred this distinction because its scripted model returned
candidates already carrying token kinds and integer scores. ENGINE-1
returns a typed \texttt{Logits} vector. The selector scans it. The
runtime uses the model/tokenizer contract to recognize EOS and otherwise
asks the tokenizer for output bytes.

\begin{quote}
\textbf{FIRST PRINCIPLE} The model produces logits. Token selection is a
separate engine policy.
\end{quote}

\Appref{sampling} owns softmax, numerical stability, temperature, randomness,
seeds, top-k, top-p, and the complete feedback loop. We will not smuggle
those topics into \texttt{forward} merely to make this chapter appear
more complete.

\section{The complete hand-calculated forward
pass}\label{the-complete-hand-calculated-forward-pass}

Our vocabulary has four identities:

\begin{verbatim}
0 <eos>
1 I
2 like
3 Rust
\end{verbatim}

The hidden dimension is three. Therefore \(V_{\mathrm{vocab}}=4\) and
\(D=3\). The parameters are:

\begin{verbatim}
E = [
  [ 0.0,  0.0,  0.0],
  [ 0.0,  1.0,  0.0],
  [ 1.0, -0.5,  2.0],
  [-1.0,  0.0,  0.0],
]

W = [
  [-0.5,  0.4,  0.1],
  [ 0.2,  0.2,  0.0],
  [ 0.3,  0.2,  0.1],
  [ 1.0, -0.4,  0.25],
]

b = [-0.2, 0.0, 0.0, 0.5]
\end{verbatim}

These values are supplied directly. In a real neural language model,
training would normally produce parameters by optimizing an objective
over data. This book is about inference, so we hold training outside the
boundary. ``Learned parameters'' names their usual origin; it does not
mean ENGINE-1 trains them.

Input \texttt{like} has ID 2. First select the embedding row:

\begin{equation}
\mathbf{h}=\mathbf{E}_{2,:}=[1.0,-0.5,2.0].
\end{equation}

Now calculate every candidate.

All four candidate rows follow the same contract:

\begin{equation}
\begin{aligned}
z_0&=-0.2+(-0.5)(1.0)+(0.4)(-0.5)+(0.1)(2.0)=-0.7,\\
z_1&= 0.0+(0.2)(1.0)+(0.2)(-0.5)+(0.0)(2.0)= 0.1,\\
z_2&= 0.0+(0.3)(1.0)+(0.2)(-0.5)+(0.1)(2.0)= 0.4,\\
z_3&= 0.5+(1.0)(1.0)+(-0.4)(-0.5)+(0.25)(2.0)=2.2.
\end{aligned}
\end{equation}

Therefore the full oracle is

\begin{equation}
\mathbf{z}=[-0.7,0.1,0.4,2.2].
\end{equation}

Notice what the oracle does not say. It does not say merely ``the answer
is token 3.'' Many wrong vectors can share the same largest index. A
transposed matrix, missing bias, wrong embedding row, or misaligned
vocabulary can still produce argmax 3 by accident.

\begin{quote}
\textbf{FIRST PRINCIPLE} A correct selected token does not prove that
the model produced correct logits.
\end{quote}

\section{An independent numerical
oracle}\label{an-independent-numerical-oracle}

The Rust model and its test could share the same indexing bug. Rewriting
the same function under a different name would not provide much
independence. The repository therefore includes a plain-Python oracle in
\texttt{\seqsplit{code/reference/python/chapter03\_oracle.py}}.

The script defines \texttt{E}, \texttt{W}, and \texttt{b} separately,
selects the row, implements the equations with its own loops, and
asserts the full expected vector. It does not import ENGINE-1, call the
Rust binary, or use NumPy. Run it with:

\begin{verbatim}
python3 code/reference/python/chapter03_oracle.py
\end{verbatim}

Expected output includes:

\begin{verbatim}
input=2:like
hidden=[1.0, -0.5, 2.0]
logits=[-0.7, 0.1, 0.4, 2.2]
oracle=PASS
\end{verbatim}

The Rust differential uses a small absolute-plus-relative tolerance:

\begin{equation}
|a-r|\le \epsilon_{\mathrm{abs}}+
\epsilon_{\mathrm{rel}}|r|,
\end{equation}

Why not exact equality? Several fixture values happen to have exact
binary representations, and this short reduction is deterministic on the
current scalar path. But floating-point arithmetic generally rounds, and
later execution providers may change reduction order. An absolute
tolerance covers small values near zero; a relative tolerance scales
with larger expected values. The tolerance must be tight enough to catch
errors and documented instead of chosen after seeing a failure.

The model also rejects non-finite parameters and non-finite output
logits. \texttt{NaN} is especially dangerous because comparisons with it
are false; a naive argmax scan can silently keep an unrelated candidate.
Some future operations may deliberately use infinities to represent
masks. ENGINE-1 has no mask or logits processor, so \texttt{NaN},
positive infinity, and negative infinity indicate invalid numerical
state at this boundary.

\section{A tensor, before a tensor
framework}\label{a-tensor-before-a-tensor-framework}

\Appref{tensors} will treat tensors systematically. We need a smaller working
definition now. A \textbf{tensor} is numerical data interpreted with
metadata and contracts such as:

\begin{itemize}
\tightlist
\item
  shape;
\item
  element type, or \textbf{dtype};
\item
  physical data;
\item
  layout and strides;
\item
  owner and lifetime;
\item
  device or location.
\end{itemize}

ENGINE-1 uses only one location, the CPU, and one storage type,
\texttt{f32}. All arrays are contiguous. Matrices are row-major. There
is no dynamic rank, broadcasting, autograd, GPU abstraction, or general
stride machinery.

The canonical
\href{https://github.com/hermonai/inside-the-llm-engine/blob/astra-visual-rewrite/diagrams/model/tiny-model-tensor-shapes.txt}{shape
table} is:

{\def\LTcaptype{none} % do not increment counter
\begin{longtable}[]{@{}lllll@{}}
\toprule\noalign{}
Name & Symbol & Shape & Lifetime & Inference-mutable? \\
\midrule\noalign{}
\endhead
\bottomrule\noalign{}
\endlastfoot
Embedding & \texttt{E} & \texttt{{[}V,D{]}} & model & no \\
Projection & \texttt{W} & \texttt{{[}V,D{]}} & model & no \\
Bias & \texttt{b} & \texttt{{[}V{]}} & model & no \\
Input token & \texttt{x} & scalar & request/step & new \\
Hidden & \texttt{h} & \texttt{{[}D{]}} & forward & new \\
Logits & \texttt{z} & \texttt{{[}V{]}} & forward/result & new \\
\end{longtable}
}

``Mutable?'' in this table describes inference semantics. A loader may
fill an allocation during construction, and Rust's
\texttt{Vec\textless{}f32\textgreater{}} is mechanically capable of
mutation. After construction, ordinary forward inference does not change
parameters. \texttt{\seqsplit{TinyLanguageModel::forward}} takes
\texttt{\&self}, making that rule visible in the API.

\begin{quote}
\textbf{FIRST PRINCIPLE} Every tensor needs an explicit shape, dtype,
layout, owner, lifetime, and access pattern. Shape is an executable
contract, not decoration.
\end{quote}

\section{Logical shape and physical row-major
memory}\label{logical-shape-and-physical-row-major-memory}

\texttt{W} has logical shape \texttt{{[}V,D{]}}. For \texttt{V=4},
\texttt{D=3}, picture four rows and three columns:

\begin{verbatim}
W[0,0] W[0,1] W[0,2]
W[1,0] W[1,1] W[1,2]
W[2,0] W[2,1] W[2,2]
W[3,0] W[3,1] W[3,2]
\end{verbatim}

The contiguous physical vector is:

\begin{verbatim}
[W00, W01, W02, W10, W11, W12, W20, W21, W22, W30, W31, W32]
\end{verbatim}

The offset for logical element \texttt{(row,\ column)} is:

\begin{verbatim}
offset = row * D + column
\end{verbatim}

This convention is not forced by the equation. We could store columns
contiguously, pad rows, tile blocks, quantize groups, or use
provider-specific layouts. Those choices would change access and kernels
without changing the model's mathematical meaning. ENGINE-1 chooses
row-major because one output candidate consumes one complete row, so the
inner loop reads contiguous weights.

Matrix notation alone cannot tell us whether bytes are contiguous,
transposed, padded, packed, resident on a device, or available through a
view. An engine must connect logical indices to physical addresses.

\section{Parameters and activations are different
memory}\label{parameters-and-activations-are-different-memory}

The embedding, projection, and bias are \textbf{parameters}: persistent
values that define model behavior. The selected hidden row and logits
are \textbf{activations}: values produced while executing the model for
a particular input.

The
\href{https://github.com/hermonai/inside-the-llm-engine/blob/astra-visual-rewrite/diagrams/model/parameters-vs-activations.txt}{ownership
diagram} shows why the distinction matters:

Figure~\ref{fig:parameter-owners} separates the model owner from request-local results. Read-only sharing never creates a shared mutable activation.

ENGINE-1 does not implement concurrency, but it chooses an ownership
model that can support it later. Multiple requests may read the same
immutable parameters. Each forward call creates its own hidden and logit
vectors. If activations were accidental mutable fields on the shared
model, concurrent requests could overwrite one another even though the
weights were correct.

This distinction will recur in memory planning, GPU execution, batching,
and the KV cache. Parameters often live as long as the loaded model.
Activations may live for one operation, one token step, one request, or
longer if the model needs persistent state.

The
\href{https://github.com/hermonai/inside-the-llm-engine/blob/astra-visual-rewrite/diagrams/model/parameters-vs-activations.txt}{memory
diagram} is small now:

\begin{equation}M_{\mathrm{parameters}}=4(12+12+4)=112\ \mathrm{bytes},\qquad M_{\mathrm{forward\ payload}}=4(3+4)=28\ \mathrm{bytes}.\end{equation}

No ``hidden reasoning'' is stored here. The adjective \emph{hidden}
distinguishes an internal representation from observable input and
output. It does not imply a human-readable chain of thought.

\section{Follow the token, byte, and
owner}\label{follow-the-token-byte-and-owner}

The
\href{https://github.com/hermonai/inside-the-llm-engine/blob/astra-visual-rewrite/diagrams/model/token-id-to-logits.txt}{complete
Appendix A.2 path} extends our first recurring journey:

\EngineFigure{model-path}{A token selects an embedding; the activation and independent projection parameters produce one score per vocabulary row.}{fig:model-path}

Follow one parameter byte conceptually. It begins as a fixture literal.
Model construction places it in a contiguous
\texttt{Vec\textless{}f32\textgreater{}}. A forward call calculates its
row-major offset, the CPU loads it, multiplies it by one hidden value,
and adds the product to an \texttt{f32} accumulator. Later chapters will
replace source literals with model-file bytes and may replace scalar
loads with packed SIMD or device operations. The ownership and
interpretation trail must remain auditable.

Follow the owner. \texttt{\seqsplit{TinyLanguageModel}} owns \texttt{E},
\texttt{W}, \texttt{b}, \texttt{V}, and \texttt{D}. A \texttt{Request}
owns its input token sequence and generation budget. The forward result
owns \texttt{h} and typed \texttt{Logits}. The runtime owns generated
token history, the UTF-8 buffer, ordered stream emissions, timings, and
the only terminal transition. The selector borrows logits and returns an
identity; it owns neither parameters nor activations.

\section{Count the parameters before counting
performance}\label{count-the-parameters-before-counting-performance}

The embedding contains \(V_{\mathrm{vocab}}D\) parameters. The untied
output projection contains another \(V_{\mathrm{vocab}}D\). The bias
contains \(V_{\mathrm{vocab}}\). Therefore:

\begin{equation}
N_{\mathrm{param}}
=V_{\mathrm{vocab}}D+V_{\mathrm{vocab}}D+V_{\mathrm{vocab}}
=2V_{\mathrm{vocab}}D+V_{\mathrm{vocab}}.
\end{equation}

With \texttt{f32}, each parameter occupies four bytes:

\begin{equation}
B_{\mathrm{param}}
=4\left(2V_{\mathrm{vocab}}D+V_{\mathrm{vocab}}\right)
\quad[\mathrm{bytes}].
\end{equation}

For ENGINE-1, \(V_{\mathrm{vocab}}=4\) and \(D=3\):

\begin{equation}
N_{\mathrm{param}}=2(4)(3)+4=28,
\qquad B_{\mathrm{param}}=28(4)=112\ \mathrm{bytes}.
\end{equation}

This count excludes \texttt{Vec} metadata, allocator bookkeeping,
dimensions, code, and forward activations. It describes parameter
payload only. The hidden activation uses \texttt{D*4\ =\ 12} bytes of
\texttt{f32} payload; the logits use \texttt{V*4\ =\ 16}.

The same formulas produce a useful scale table:

{\def\LTcaptype{none} % do not increment counter
\begin{longtable}[]{@{}
  >{\raggedleft\arraybackslash}p{(\linewidth - 10\tabcolsep) * \real{0.1667}}
  >{\raggedleft\arraybackslash}p{(\linewidth - 10\tabcolsep) * \real{0.1667}}
  >{\raggedleft\arraybackslash}p{(\linewidth - 10\tabcolsep) * \real{0.1667}}
  >{\raggedleft\arraybackslash}p{(\linewidth - 10\tabcolsep) * \real{0.1667}}
  >{\raggedleft\arraybackslash}p{(\linewidth - 10\tabcolsep) * \real{0.1667}}
  >{\raggedleft\arraybackslash}p{(\linewidth - 10\tabcolsep) * \real{0.1667}}@{}}
\toprule\noalign{}
\begin{minipage}[b]{\linewidth}\raggedleft
\texttt{V}
\end{minipage} & \begin{minipage}[b]{\linewidth}\raggedleft
\texttt{D}
\end{minipage} & \begin{minipage}[b]{\linewidth}\raggedleft
\texttt{E} parameters
\end{minipage} & \begin{minipage}[b]{\linewidth}\raggedleft
\texttt{W} parameters
\end{minipage} & \begin{minipage}[b]{\linewidth}\raggedleft
Total \texttt{2VD+V}
\end{minipage} & \begin{minipage}[b]{\linewidth}\raggedleft
\texttt{f32} bytes
\end{minipage} \\
\midrule\noalign{}
\endhead
\bottomrule\noalign{}
\endlastfoot
4 & 3 & 12 & 12 & 28 & 112 \\
100 & 16 & 1,600 & 1,600 & 3,300 & 13,200 \\
50,000 & 4,096 & 204,800,000 & 204,800,000 & 409,650,000 &
1,638,600,000 \\
\end{longtable}
}

The formula scales quickly. Consider a hypothetical vocabulary of 50,000
and hidden dimension 4,096 under this same simplified, untied
architecture:

\begin{equation}
\begin{aligned}
V_{\mathrm{vocab}}D&=50{,}000(4{,}096)=204{,}800{,}000,\\
N_{\mathrm{param}}&=409{,}650{,}000,\\
B_{\mathrm{param}}&=1{,}638{,}600{,}000\ \mathrm{bytes}
\approx1.526\ \mathrm{GiB}.
\end{aligned}
\end{equation}

That is not the size of a complete real language model and not a
measurement. It isolates vocabulary embedding and projection under
stated assumptions. It does explain why output heads are not negligible
and why real designs care about lower-precision storage, quantization,
and optimized multiplication.

Some models use \textbf{weight tying}: the input embedding and output
projection reuse the same parameter storage, with the orientation
interpreted according to the chosen convention. That can eliminate one
\texttt{V*D} parameter collection. It is an architectural option, not a
universal law. ENGINE-1 leaves \texttt{E} and \texttt{W} separate
because changing an input representation and changing a candidate
scoring row should be visibly different experiments.

\section{Count the work and the bytes
read}\label{count-the-work-and-the-bytes-read}

Embedding lookup performs little arithmetic. It validates \texttt{x},
computes an offset, and reads one row: \texttt{D} parameter values. It
does not scan all \texttt{V} rows.

Output projection is different. For each of \texttt{V} output rows, it
performs \texttt{D} multiplications and approximately \texttt{D}
additions including bias accumulation. The work grows as

\begin{equation}
\Theta\!\left(V_{\mathrm{vocab}}D\right)
\end{equation}

for this exact dense projection: every output row and hidden component
is visited once.

For one forward call, the scalar kernel reads essentially all
\texttt{V*D} projection weights and \texttt{V} biases and writes
\texttt{V} logits. It reuses the small hidden vector across all rows.
Even in this tiny model, two distinct systems patterns appear:

\begin{itemize}
\tightlist
\item
  embedding lookup: sparse row access, little arithmetic;
\item
  output projection: broad weight access plus many dot-product
  operations.
\end{itemize}

Inference performance depends on computation and data movement. A
processor cannot multiply a weight it has not obtained through its
memory hierarchy. Later chapters will ask whether a kernel is limited by
arithmetic, weight bandwidth, cache reuse, conversion, launch cost,
synchronization, or another resource. \Appref{smallest-model} establishes the
accounting habit without pretending that a 4-by-3 loop predicts a
production workload.

The repository includes
\texttt{\seqsplit{code/experiments/chapter-03-projection-scaling.py}}.
It varies modest \texttt{V} and \texttt{D}, reports parameter payload
and median time for a plain-Python scalar loop, and prints a checksum.
Its result record is in
\texttt{\seqsplit{research/benchmarks/chapter-03-projection-scaling.md}}.

\begin{quote}
\textbf{PERFORMANCE LAB} The scaling probe is pedagogical. It shows that
larger shapes create more stored values and loop iterations in that
harness. Python interpreter timing does not predict Rust, BLAS, SIMD,
GPU, or real-model throughput, and it must not be compared with Hermon
or llama.cpp.
\end{quote}

\section{Model semantics are not one
implementation}\label{model-semantics-are-not-one-implementation}

The equations say what the model computes:

\begin{verbatim}
h = E[x]
z = W h + b
\end{verbatim}

The scalar Rust loops say how ENGINE-1 executes those semantics today.
The
\href{https://github.com/hermonai/inside-the-llm-engine/blob/astra-visual-rewrite/diagrams/model/semantics-vs-execution.txt}{semantics/execution
diagram} keeps that line visible:

\begin{quote}
\begin{enumerate}
\item MODEL SEMANTICS
\item TokenId  \(\longrightarrow\)  embedding  \(\longrightarrow\)  hidden  \(\longrightarrow\)  projection  \(\longrightarrow\)  logits
\item EXECUTION IMPLEMENTATION
\item today: scalar row-major Rust loops
\item later: CPU SIMD, Metal, CUDA, other native kernels
\end{enumerate}
\end{quote}

An optimized implementation may load several weights into vector
registers, tile a matrix, use a quantized packed layout, fuse
operations, or dispatch to a device. It can allocate scratch
differently. None of those choices may change which token owns a row,
transpose the mathematical convention accidentally, drop the bias, or
produce errors outside the declared numerical tolerance.

\begin{quote}
\textbf{FIRST PRINCIPLE} An engine may change execution while preserving
model semantics. The oracle defines what must remain invariant.
\end{quote}

This separation is the foundation of later differential testing. The
clarity path need not be fast. The fast path must agree with an
independent clarity path before its timing is meaningful.

\section{BUILD IT: ENGINE-1}\label{build-it-engine-1}

ENGINE-1 evolves the existing Rust crate instead of maintaining parallel
copies of an entire runtime. Git history preserves the exact ENGINE-0
milestone. The current code keeps \Appref{tokens}'s tokenizer trait,
chat/template types, request, streaming events, UTF-8 framer,
cancellation contract, timings, and terminal lifecycle. It replaces the
historical scripted center; the current checkout integrates the later tensor
and sampling implementations while preserving this numerical fixture.

\subsection{The model structure}\label{the-model-structure}

\texttt{\seqsplit{TinyLanguageModel}} owns:

\begin{verbatim}
pub struct TinyLanguageModel {
    vocab_size: usize,
    hidden_dim: usize,
    embedding: OwnedTensor,
    output_weight: OwnedTensor,
    output_bias: Vec<f32>,
}
\end{verbatim}

Read \texttt{OwnedTensor} here as an owned contiguous array with checked
shape metadata; \Appref{tensors} derives its implementation. The historical
ENGINE-1 milestone used bare vectors, and the current checkout preserves
its values through this explicit storage wrapper. Field names carry the
semantic rank and role. Construction checks every length before execution.

The model trait accepts the full history:

\begin{verbatim}
pub trait Model {
    fn vocabulary_size(&self) -> usize;
    fn forward(&self, input: &[TokenId])
        -> Result<ForwardPass, ModelError>;
}
\end{verbatim}

Why accept a sequence if ENGINE-1 uses only its last element? Passing
only \texttt{last\_token} would be locally narrower, but it would make
the runtime decide which history information matters to the model. The
full borrowed slice keeps next-token history at the model boundary,
introduces no input allocation, and makes the limitation testable.
\texttt{forward} calls \texttt{input.last()} and says so in code. Future
models may consume more positions without changing what a request means.

This is a documented teaching choice, not a claim that every production
model API should use the same signature. Batched engines need shapes,
positions, sequence identifiers, cache handles, and output-row requests
that this method does not express.

\subsection{A typed forward result}\label{a-typed-forward-result}

\texttt{ForwardPass} contains:

\begin{verbatim}
input_token: TokenId
hidden:      Vec<f32>
logits:      Logits
\end{verbatim}

Keeping the hidden activation in the teaching result makes trace and
oracle inspection straightforward. A production API might retain it only
in debug mode or place scratch in a reusable context. \texttt{Logits}
wraps its vector and exposes a borrowed slice. This prevents unrelated
code from treating every \texttt{Vec\textless{}f32\textgreater{}} as
interchangeable.

The model owns \texttt{Logits} construction and guarantees finite
values. It returns a typed \texttt{ModelError} for an empty sequence or
out-of-range final token. The selector receives only \texttt{\&Logits},
so it cannot reach into parameters or alter the hidden activation.

\subsection{Construction is a correctness
boundary}\label{construction-is-a-correctness-boundary}

\texttt{\seqsplit{TinyLanguageModel::try\_new}} rejects:

\begin{itemize}
\tightlist
\item
  zero vocabulary size;
\item
  zero hidden dimension;
\item
  overflow in \texttt{V*D};
\item
  an embedding length other than \texttt{V*D};
\item
  a projection length other than \texttt{V*D};
\item
  a bias length other than \texttt{V};
\item
  any non-finite parameter.
\end{itemize}

These are load-time properties. Paying once during construction lets the
hot inner loop index validated storage without rechecking every row
length. Forward still checks the request-dependent token ID.

Shape bugs are not cosmetic. If \texttt{W} has 11 values instead of 12,
choosing to truncate would omit a term, choosing to pad would invent a
parameter, and indexing blindly could panic. None preserves the declared
model.

\begin{quote}
\textbf{PROVE IT} Model shape validation is inference correctness.
Malformed parameter counts and vocabulary mismatches must fail before
plausible output can hide them.
\end{quote}

\subsection{Vocabulary coupling becomes
numerical}\label{vocabulary-coupling-becomes-numerical}

\Appref{tokens} established that a tokenizer revision and chat template belong
to a model contract. \Appref{smallest-model} adds a numeric equality:

\begin{verbatim}
model vocabulary size
    == tokenizer vocabulary size
    == ModelContract vocabulary size
\end{verbatim}

Embedding row \texttt{i}, logit position \texttt{i}, tokenizer identity
\texttt{i}, and EOS recognition must refer to the same token. A vector
of length four cannot safely consume \Appref{tokens}'s sparse byte-BPE IDs,
which extend through 1007.

ENGINE-1 therefore has a deliberately tiny paired tokenizer:

\begin{verbatim}
0 <eos>   special, no ordinary output bytes
1 I       bytes "I"
2 like    bytes " like"
3 Rust    bytes " Rust"
\end{verbatim}

Leading spaces belong to two pieces, allowing \texttt{I\ like\ Rust} to
round-trip. The tokenizer rejects all other text with a byte-offset
error. That narrow domain is honest: four embeddings cannot represent a
thousand-ID vocabulary. \Appref{tokens}'s byte tokenizer and BPE remain in the
repository for their own labs; they are not silently remapped into
ENGINE-1.

\texttt{\seqsplit{Runtime::try\_new}} validates the model, tokenizer,
and contract sizes and requires an EOS identity. It cannot admit a
request when the model's fourth logit would mean one token to the model
and another to the decoder.

\subsection{The inherited lifecycle}\label{the-inherited-lifecycle}

For prompt \texttt{I\ like}, the tokenizer yields \texttt{{[}1,2{]}}.
The first forward uses the last token 2 and produces the oracle logits.
Temporary argmax selects token 3, which decodes to bytes for
\texttt{Rust}. The runtime emits a token event and then a valid text
event.

The generated token is appended to history. The next forward therefore
ends in token 3. With the fixture's \texttt{Rust} embedding
\texttt{{[}-1.0,0.0,0.0{]}}, the logits are:

\begin{verbatim}
[0.3, -0.2, -0.3, -0.5]
\end{verbatim}

EOS at position 0 is the unique maximum. The runtime recognizes it as a
terminal identity, does not ask for ordinary text bytes, checks that the
UTF-8 buffer is complete, and emits exactly one completed terminal
event.

This two-step behavior is not a disguised step table. Both logit vectors
come from the same immutable \texttt{E}, \texttt{W}, and \texttt{b};
only the input embedding row changes.

The lifecycle retains its earlier failure properties:

\begin{itemize}
\tightlist
\item
  empty input or zero generation budget fails before admission;
\item
  cancellation stops before another model call or token emission;
\item
  model, tokenizer, and UTF-8 failures produce one failed terminal;
\item
  EOS produces no text token event;
\item
  maximum-token termination checks pending UTF-8 bytes;
\item
  no event or trace can occur after terminal;
\item
  one request cannot both complete and fail.
\end{itemize}

\Appref{sampling} will explain the loop as autoregressive generation. Here it
remains minimal infrastructure needed to prove that the new logit
boundary integrates with established ownership.

\subsection{An educational numerical
trace}\label{an-educational-numerical-trace}

Run:

\begin{verbatim}
cd code/mini-engine
cargo run -p engine0 -- --trace 'I like'
\end{verbatim}

The trace names admission and execution, then includes the input token,
hidden shape and values, logit shape and values, selected token, decoded
byte count, streamed text, and terminal reason. ENGINE-1 vectors are
small enough to show completely. This must not become a habit of dumping
future tensors with millions of elements. The trace already names shapes
separately from values; later milestones can retain shapes, dtypes,
selected values, and statistics while bounding payload size.

\section{Test the arithmetic, not the
story}\label{test-the-arithmetic-not-the-story}

The Rust suite checks more than one happy result:

\begin{itemize}
\tightlist
\item
  valid construction and exact parameter/byte counts;
\item
  zero \texttt{V}, zero \texttt{D}, and multiplication overflow;
\item
  embedding, projection, and bias count mismatches;
\item
  non-finite parameters;
\item
  empty model input and an ID outside \texttt{0..V};
\item
  exact embedding-row selection;
\item
  the complete oracle logit vector within tolerance;
\item
  bias behavior with a zero embedding;
\item
  negative and zero arithmetic;
\item
  repeatability of the complete forward result;
\item
  identical logits for different histories with the same final token;
\item
  expected different logits for different embedding rows;
\item
  EOS as an ordinary position in the output vector;
\item
  one changed projection weight affecting only its row's logit;
\item
  model/tokenizer/contract vocabulary mismatch;
\item
  runtime cancellation, failure, EOS, max-token, UTF-8, event ordering,
  and exactly-one-terminal behavior.
\end{itemize}

Changing one projection weight is a particularly useful causality test.
For input \texttt{like}, change \texttt{W{[}3,0{]}} from 1.0 to 1.5.
Since \texttt{h\_0=1.0}, only \texttt{z\_3} increases, by 0.5:

\begin{verbatim}
before: [-0.7, 0.1, 0.4, 2.2]
after:  [-0.7, 0.1, 0.4, 2.7]
\end{verbatim}

The embedding does not change. Rows 0 through 2 do not read
\texttt{W{[}3,0{]}}. By contrast, changing one embedding component can
affect every logit because all projection rows consume the shared hidden
component.

Labs 5--8 turn these invariants into exercises: calculate the full pass,
change one weight, prove the context limitation, and break the shape
deliberately.

\section{Inside Hermon: the same boundary, industrial
machinery}\label{inside-hermon-the-same-boundary-industrial-machinery}

\begin{quote}
\textbf{INSIDE HERMON --- CURRENT} At inspected commit \texttt{472a44c},
Hermon's default runtime is the continuous- batched path, and real-model
numerical execution is delegated through \texttt{hermon-llamacpp} to
pinned llama.cpp commit \texttt{389ff61}.
\end{quote}

Hermon's default path is not an embedding-plus-one-projection model. It
runs a real model graph with layers, context state, batching, and
hardware backends. The useful comparison is the input/output contract.

\texttt{\seqsplit{crates/hermon-runtime/src/dispatch.rs}} selects
\texttt{\seqsplit{RuntimeMode::Batched}} when
\texttt{\seqsplit{HERMON\_RUNTIME\_MODE}} is absent.
\texttt{\seqsplit{crates/hermon-runtime/src/batched.rs}} owns request
scheduling and assembles token positions from multiple active sequences
into a \texttt{Batch}. Each entry can request that its output logits be
kept. The worker calls one context decode and samples the requested
logit row for each sequence.

The simpler blocking path in
\texttt{\seqsplit{crates/hermon-engine/src/engine.rs}} makes the journey
easy to follow: clear per-request context state, tokenize the prompt,
construct a sampler, call context generation, turn returned token IDs
into pieces, frame UTF-8, and emit text. Inside
\texttt{hermon-llamacpp}, \texttt{Context::decode} passes token IDs
through a C shim to \texttt{llama\_decode}.

The pinned llama.cpp header defines the numerical output contract. After
\texttt{llama\_decode}, logits requested by the batch are stored as
rows; each row has \texttt{n\_vocab} columns.
\texttt{\seqsplit{llama\_get\_logits\_ith}} retrieves one output row.
The context implementation synchronizes pending computation before
returning that result.

Hermon therefore owns substantial systems machinery around the current
forward call---model resolution, runtime choice, batch construction,
scheduling, streaming, and lifecycle---while llama.cpp owns the default
real-model numerical graph and its output-logit storage. A sampler
consumes those logits through the llama.cpp context. ENGINE-1 exposes
\texttt{Logits} directly because the teaching goal is to make that
otherwise industrial boundary visible.

\begin{quote}
\textbf{INSIDE HERMON --- PREVIEW} Hermon's own paged GGUF forward is
present behind \texttt{\seqsplit{HERMON\_RUNTIME\_MODE=paged}} and
\texttt{\seqsplit{HERMON\_PAGED\_GGUF=1}}, with a currently stated
greedy CPU limit. Its real-model differential test compares next-token
behavior with the pinned llama.cpp graph. The gate means this path is
not the default CURRENT numerical engine.
\end{quote}

Hermon also contains tensor, backend, paged-KV, and native-kernel
components. Their existence is \textbf{LIBRARY} evidence unless the
request path selects them. This distinction prevents a common
source-reading error: finding a sophisticated kernel and describing
every production request as if it executes there.

ENGINE-1 resembles none of Hermon's production model architecture. It
does establish the abstract contract both systems need:

\begin{quote}
\begin{enumerate}
\item token identities  \(\longrightarrow\)  model execution  \(\longrightarrow\)  one vocabulary score per output row
\end{enumerate}
\end{quote}

\section{The fatal context
limitation}\label{the-fatal-context-limitation}

Return to the approximation:

\begin{verbatim}
P(next | history) is approximated by P(next | final token)
\end{verbatim}

Consider two histories:

\begin{quote}
\begin{enumerate}
\item History A: I  \(\longrightarrow\)  like     \(\longrightarrow\)  Rust
\item History B: I  \(\longrightarrow\)  dislike  \(\longrightarrow\)  Rust
\end{enumerate}
\end{quote}

Both end in the same \texttt{TokenId} for \texttt{Rust}. ENGINE-1
selects the same \texttt{E{[}Rust{]}}, produces the same hidden vector,
and calculates the same logits. The
\href{https://github.com/hermonai/inside-the-llm-engine/blob/astra-visual-rewrite/diagrams/model/context-limitation.txt}{context
limitation diagram} shows information disappearing at the model
boundary.

This is not merely theoretical. The test
\texttt{\seqsplit{same\_last\_token\_produces\_same\_logits\_despite\_different\_history}}
passes two different-length sequences with different prefixes and the
same final token. It asserts equality of both hidden vectors and the
complete logit vectors. A control with a different final token produces
a different embedding and the expected different logits.

No choice of argmax, softmax, temperature, or sampling can recover
history the model never used. Selection policies operate on the logits
they receive. The missing capability must arise before logits:

\begin{verbatim}
We need a mechanism that makes the representation at the current position
depend on context.
\end{verbatim}

Then stop. Appendices A.4--A.8 build that mechanism progressively. Jumping
directly to attention here would skip the tensor, multiplication,
normalization, projection, position, and ownership foundations needed to
implement it correctly.

\section{Common mistakes}\label{common-mistakes-2}

\textbf{Treating token IDs as ordinal measurements.} IDs are vocabulary
identities. Only their mapping to parameter rows supplies numerical
semantics.

\textbf{Calling embedding lookup a required dense multiplication.}
One-hot multiplication is an equivalent equation, not ENGINE-1's
physical execution. The implementation reads one row.

\textbf{Giving hidden dimensions human definitions.} A hidden vector is
a learned activation consumed by later computation. Individual
coordinates need not have isolated, stable meanings.

\textbf{Calling the hidden vector a probability distribution.} Its
values can be arbitrary finite model activations. They need not be
positive or sum to one.

\textbf{Calling logits probabilities.} Logits are unnormalized scores.
Negative values and totals other than one are normal.

\textbf{Assuming the largest logit has probability one.} Argmax returns
an index. It does not define probability or uncertainty.

\textbf{Calling argmax the model.} The model produces the score vector.
Argmax is one selection policy consuming it.

\textbf{Putting selection inside \texttt{forward}.} Doing so prevents
independent testing of logits and entangles immutable model semantics
with generation policy.

\textbf{Confusing parameters and activations.} \texttt{E}, \texttt{W},
and \texttt{b} persist across requests. \texttt{h} and \texttt{z} belong
to a forward execution.

\textbf{Mutating weights during inference.} Ordinary forward execution
reads the loaded model. Training or adaptation is a separate operation
and ownership contract.

\textbf{Testing only the selected token.} Wrong logit values can
preserve the same argmax. Compare the whole vector with an independent
oracle.

\textbf{Passing anonymous vectors without shape meaning.}
\texttt{Vec\textless{}f32\textgreater{}} says neither whether values are
hidden activations or logits nor which vocabulary owns their positions.
Lightweight types and construction checks carry the contract.

\textbf{Treating shape as prose.} A dimension mismatch changes or
invalidates the computation. Reject it before indexing.

\textbf{Assuming matrix notation specifies memory.}
\texttt{z\ =\ W\ h\ +\ b} does not tell us row-major versus column-major
storage, padding, dtype, location, or ownership.

\section{Exercises}\label{exercises-2}

\subsection{CHECK}\label{check-2}

\begin{enumerate}
\def\labelenumi{\arabic{enumi}.}
\tightlist
\item
  For \texttt{V=8}, \texttt{D=5}, give the shapes of \texttt{E},
  \texttt{W}, \texttt{b}, \texttt{h}, and \texttt{z}.
\item
  Calculate the untied parameter count and \texttt{f32} payload bytes.
\item
  If input token 6 is valid, how many embedding values does lookup read?
\item
  Explain why logit position 6 and embedding row 6 must share vocabulary
  identity even though they serve different computations.
\end{enumerate}

\subsection{BUILD}\label{build-2}

Complete
\href{https://github.com/hermonai/inside-the-llm-engine/blob/astra-visual-rewrite/labs/lab-05-forward-pass-by-hand.md}{Lab
5}. Write all four dot products before running either oracle. Then add a
hand oracle for input \texttt{Rust} and verify
\texttt{\seqsplit{[0.3,-0.2,-0.3,-0.5]}}.

\subsection{BREAK}\label{break-2}

Complete
\href{https://github.com/hermonai/inside-the-llm-engine/blob/astra-visual-rewrite/labs/lab-08-break-the-shape.md}{Lab
8}. Try malformed lengths, zero dimensions, a non-finite parameter, an
out-of-range ID, and a vocabulary mismatch. Identify which checks occur
once at model construction and which remain request-dependent.

\subsection{EXTEND}\label{extend-2}

Complete
\href{https://github.com/hermonai/inside-the-llm-engine/blob/astra-visual-rewrite/labs/lab-06-change-one-weight.md}{Lab
6} and
\href{https://github.com/hermonai/inside-the-llm-engine/blob/astra-visual-rewrite/labs/lab-07-same-last-token-same-output.md}{Lab
7}. Then add a second small valid fixture with \texttt{V=3},
\texttt{D=2}. Keep all values hand-computable and write the expected
full vector independently before implementing the Rust case.

\section{Summary}\label{summary-2}

A token ID contains identity, not geometry. An embedding matrix gives
that identity a \texttt{D}-component numerical representation by
selecting one row. An output projection compares the hidden vector with
one row per vocabulary candidate and adds a bias:

\begin{verbatim}
h = E[x]
z = W h + b
\end{verbatim}

ENGINE-1 makes each shape, dtype, layout, owner, lifetime, and scalar
operation explicit. Its \texttt{V=4}, \texttt{D=3} fixture has 28
parameters occupying 112 bytes of \texttt{f32} payload. For input
\texttt{like}, it produces the independently verified vector:

\begin{verbatim}
[-0.7, 0.1, 0.4, 2.2]
\end{verbatim}

The fake candidate table is gone from the current model path. Model
execution returns typed finite logits; a separate temporary selector
chooses an ID. The request lifecycle, EOS handling, byte decoding,
strict UTF-8 framing, cancellation, failure, and exactly-one-terminal
invariant remain executable.

The model is intentionally inadequate. It uses only the final token, so
any two histories ending in that token produce identical logits. That
failure creates the need for context-dependent representations without
hiding the first numerical inference step behind Transformer complexity.

\section{What comes next}\label{what-comes-next}

We now have arbitrary-looking but correct scores:

\begin{verbatim}
[-0.7, 0.1, 0.4, 2.2]
\end{verbatim}

How should an engine turn them into the next token? \Appref{sampling}---\textbf{Logits, Sampling, and the Autoregressive Loop}---will derive
numerically stable softmax, distinguish greedy and categorical
selection, introduce temperature, seeds, top-k and top-p, append the
selected identity to history, execute the model again, and reconcile EOS
and token limits with terminal ownership.

The separation established here must survive:

\begin{quote}
\begin{enumerate}
\item model semantics  \(\longrightarrow\)  logits  \(\longrightarrow\)  logits processing  \(\longrightarrow\)  sampler  \(\longrightarrow\)  runtime stopping
\end{enumerate}
\end{quote}

\Appref{sampling} will not have to invent logits or guess what they mean. It can
begin at a typed, oracle-verified full vocabulary vector.

\section{References}\label{references}

\begin{itemize}
\tightlist
\item
  Yoshua Bengio, Réjean Ducharme, Pascal Vincent, and Christian Jauvin,
  \href{https://jmlr.org/papers/v3/bengio03a.html}{``A Neural
  Probabilistic Language Model''}, \emph{Journal of Machine Learning
  Research} 3, 2003.
\item
PyTorch,
  \href{https://docs.pytorch.org/docs/stable/generated/torch.nn.Embedding}{\texttt{\seqsplit{torch.nn.Embedding}}},
  official lookup-table contract, inspected 2026-09-02.
\item
PyTorch,
  \href{https://docs.pytorch.org/docs/stable/generated/torch.nn.Linear.html}{\texttt{torch.nn.Linear}},
  official affine-layer contract, inspected 2026-09-02.
\item
  Hermon source at commit
  \texttt{\seqsplit{472a44cdb511b2dae6c9569e59543db8f8350b25}},
  especially \texttt{\seqsplit{crates/hermon-engine/src/engine.rs}},
  \texttt{\seqsplit{crates/hermon-runtime/src/dispatch.rs}},
  \texttt{\seqsplit{crates/hermon-runtime/src/batched.rs}}, and
  \texttt{\seqsplit{crates/hermon-llamacpp/src/linked.rs}}.
\item
  llama.cpp source pinned by Hermon at commit
  \texttt{\seqsplit{389ff61d77b5c71cec0cf92fe4e5d01ace80b797}},
  especially \texttt{include/llama.h},
  \texttt{\seqsplit{src/llama-context.cpp}}, and
  \texttt{\seqsplit{src/llama-sampler.cpp}}.
\end{itemize}


\clearpage
\section{Worked synthesis: intervene on one parameter}
\subsection{Change one weight and predict the whole difference}
\textbf{Problem.} For $h=(1,-.5,2)^\top$, increase $W_{3,2}$ by $0.1$ and
leave every other parameter fixed. What changes?

\textbf{Solution.} Only the fourth logit changes:
\begin{equation}
\Delta z_i=\begin{cases}0.1h_2=0.2&i=3,\\0&i\ne3.\end{cases}
\end{equation}
Thus the original $z_3=2.2$ becomes $2.4$, up to F32 rounding. Greedy output
remains ID 3. A test that checks only the chosen ID would miss an incorrect
implementation that ignored the changed weight; compare the full logit vector.

\subsection{Prove the context limitation}
\textbf{Problem.} Can this model distinguish histories $[1,2]$ and $[3,2]$?

\textbf{Solution.} No. Its forward function chooses the final ID, gathers the
same embedding row and applies the same projection and bias. Both histories
produce identical logits. The runtime retains history; this model does not
use the earlier entries. Larger history allocation alone cannot create context.

\subsection{Count the persistent payload}
\textbf{Problem.} Count the bytes for $E:[4,3]$, $W:[4,3]$ and $b:[4]$ in F32.

\textbf{Solution.} $(4\cdot3+4\cdot3+4)\cdot4=112$ bytes of parameter payload.
The hidden vector adds 12 bytes and the logits add 16 bytes for one forward
result. These counts exclude vector metadata, allocator bookkeeping and history.
They are not process resident-set measurements.

